\hypertarget{cont4:polynomial-roots-under-coefficient-perturbation}{%
\subsection{Polynomial roots under coefficient
perturbation}\label{cont4:polynomial-roots-under-coefficient-perturbation}}

The argument proves uniqueness of roots for square systems and nearby-root stability
for finite systems of integer polynomials. The polynomials and their
partial derivatives are evaluated directly. No root-uniqueness or
perturbation-stability premise is assumed.

\hypertarget{cont4:notation}{%
\subsubsection{Notation}\label{cont4:notation}}

Let \(p\) be prime, \(L\ge1\), and \[
F,H\in\mathbb Z[X_1,\ldots,X_d]^e,
\qquad x\in\mathbb Z^d.
\] Write \[
A_x=DF(x)\bmod p:\mathbb F_p^d\longrightarrow\mathbb F_p^e,
\qquad B_L(x)=x+p^L\mathbb Z_p^d.
\] For an integer \(c\), the perturbed system is the actual polynomial
list \[
G=F+cH.
\] There is no restriction on the sign of \(c\), and \(c=0\) is allowed.
Divisibility conclusions give lower bounds on precision and allow zero
differences.

\hypertarget{cont4:unique-roots-for-square-systems}{%
\subsubsection{1. Unique roots for square
systems}\label{cont4:unique-roots-for-square-systems}}

If \(d=e\), \(A_x\) is surjective, and \(F(x)\equiv0\pmod{p^L}\), then
\[
\boxed{\exists!\,a\in B_L(x):F(a)=0.}
\] The stronger comparison lemma proves that any two exact roots of
\(F\) in \(B_L(x)\) coincide, without needing a separate lower-root
hypothesis.

\textbf{Proof implemented.} Map every arbitrary p-adic root into the
actual residue rings \(\mathbb Z/p^n\mathbb Z\), and choose its
canonical integer representative. Evaluation commutes with reduction, so
these representatives are genuine finite polynomial roots. The preceding
count theorem gives exactly \(p^{d-d}=1\) next-level residue lift.
Inductively any two roots agree at every precision. The p-adic residue
maps separate points, so the roots are equal. Existence is the preceding
checked tower theorem.

This extends uniqueness from the limit of a selected sequence to
\textbf{all exact roots in the prescribed ball}.

\hypertarget{cont4:coefficient-perturbations-preserve-smooth-liftability}{%
\subsubsection{2. Coefficient perturbations preserve smooth
liftability}\label{cont4:coefficient-perturbations-preserve-smooth-liftability}}

If \(p^L\mid c\), the actual identities give \[
G(x)=F(x)+cH(x),\qquad DG(x)\equiv DF(x)\pmod p.
\] Thus if \(F(x)\equiv0\pmod{p^L}\) and \(A_x\) is surjective, then \[
\exists b\in B_L(x):G(b)=0.
\] This theorem allows rectangular systems. The evaluation, derivative
and Jacobian identities are proved from the polynomial definition
\(G_i=F_i+C(c)H_i\), rather than supplied as hypotheses. Algebraic
congruence helpers also permit composite moduli; exact p-adic root
existence uses prime \(p\).

\hypertarget{cont4:nearby-roots-for-rectangular-systems}{%
\subsubsection{3. Nearby roots for rectangular
systems}\label{cont4:nearby-roots-for-rectangular-systems}}

Let \(M\ge1\) and \(p^M\mid c\). Suppose \(A_x\) is surjective and
\(a\in\mathbb Z_p^d\) is an exact root of \(F\) with
\(a\equiv x\pmod p\). Then \[
\boxed{\exists b\in\mathbb Z_p^d:\quad G(b)=0,
\qquad b\equiv a\pmod{p^M}.}
\] No uniqueness is claimed when \(d>e\).

\textbf{Proof implemented.} Truncate \(a\) to an integer vector \(y\)
modulo \(p^M\). It is an actual lower root of \(F\), and its reduced
Jacobian agrees with \(A_x\). Apply the coefficient-perturbed existence
theorem at precision \(M\), then combine the two congruences through
\(y\).

\hypertarget{cont4:full-precision-for-the-unique-square-root}{%
\subsubsection{4. Full precision for the unique root of a square
system}\label{cont4:full-precision-for-the-unique-square-root}}

Suppose \(d=e\), \(M\ge L\ge1\), \(p^M\mid c\), and \(A_x\) is
surjective. If \[
a,b\in B_L(x),\qquad F(a)=0,\qquad G(b)=0,
\] then \[
\boxed{b-a\in p^M\mathbb Z_p^d.}
\] In particular, for \(G=F+p^MH\), no p-adic digits below depth \(M\)
are changed in the selected root.

The formalization also proves a combined existence theorem. Starting
only with \(F(x)\equiv0\pmod{p^L}\), it constructs an original root
\(a\in B_L(x)\) and proves that there is exactly one \(b\in B_L(x)\)
satisfying \(G(b)=0\) and \(b\equiv a\pmod{p^M}\). The comparison
theorem shows that the precision bound holds for every root of \(G\) in
that ball, so uniqueness is not obtained merely by adding a restrictive
displacement condition.

\textbf{Proof implemented.} Section 3 produces a nearby root \(b'\).
Since \(M\ge L\), it belongs to \(B_L(x)\). The perturbed Jacobian
remains surjective, and square uniqueness identifies \(b'\) with the
arbitrary root \(b\).
